% SPDX-License-Identifier: CC0-1.0
% To the extent possible under law, the author has waived all copyright and related or
% neighboring rights to this file. See the LICENSE file (CC0 1.0 Universal).

% Supplement to zzn.tex. Every technical statement carries a comment
% "% lean: <file>:<line> <name>" on the line before it, with paths relative to proof/.
% Section numbers "Section 3.x" refer to the sections of the paper zzn.tex as built.

\documentclass[11pt]{article}
\usepackage[margin=1in]{geometry}
\usepackage{amsmath,amssymb,amsthm}
\usepackage{booktabs}
\usepackage{algorithm}
\usepackage[noend]{algpseudocode}
\usepackage{longtable}
\usepackage{array}
\usepackage{hyperref}

\newtheorem{lemma}{Lemma}
\newcommand{\lean}[1]{\texttt{#1}}
\newcommand{\lf}[1]{\nolinkurl{#1}}
\setlength{\emergencystretch}{3em}
\algrenewcommand\algorithmicrequire{\textbf{Input:}}
\algrenewcommand\algorithmicensure{\textbf{Output:}}

\title{Supplement to: Two-Colour Zig-Zag Numberlink on Rectangles:\\
A Forbidden-Pattern Characterization and a Polynomial-Time\\
Construction, Verified in Lean}
\author{Z. Zzyzek\thanks{To the extent possible under law, the author has waived all copyright and
related or neighboring rights to this supplement, under CC0 1.0 Universal
(\url{https://creativecommons.org/publicdomain/zero/1.0/}).}}
\date{Draft, November 2026}

\begin{document}
\maketitle

This supplement holds reference material for the paper. Each section supports one part of the
paper, named in its first sentence, and uses the paper's definitions and entry names (P, T1, T2,
L1 to L6, E1, B, C, D, R). Every statement is checked against the Lean development in
\lean{proof/}. The numerical values of $N(b)$ and $K(b)$ for $b=10$ and $b=16$ are computed outside
Lean, by \lean{paper/tools/kwr.py}.

\section{The plug DP encoding}\label{sec:plugdp}

This section gives the encoding of the plug DP of Section~3.1, which the window certificates of
Section~3.2 also run.

% lean: ZZN/Statement.lean:661 dpAcceptF
% lean: ZZN/Statement.lean:663 thinBF
% lean: ZZN/Statement.lean:622 cell
The instance is transposed when $h>w$, so that the grid has $a$ rows of $b\le a$ cells, where a
row is a line $x=k$. A window of side $k$ has $a=b=k$. The cells are processed in row-major order,
cell $i$ being $(\lfloor i/b\rfloor, i \bmod b)$.
% lean: ZZN/Thin/Inside.lean:432 Piece (a run that is not the whole path)
A \emph{piece} is a maximal run of a solution path inside the processed cells that is not the
whole path. Its two ends are coded as natural numbers.

% lean: ZZN/Statement.lean:558 termE
% lean: ZZN/Statement.lean:561 exitE
% lean: ZZN/WinCheck.lean:18 (module docstring: end codes)
% lean: ZZN/WinCheck.lean:130 exitCell
\begin{itemize}
\item $c<b$ is a plug down from column $c$ into the next row.
\item $b$ is a plug right, into the next cell of the row.
\item $b+1+s$ is exit slot $s$. Slot $s<a$ leaves row $s$ to the right, and slot $a+j$ leaves column
  $b-1-j$ downward. The slots are numbered down the right side and then leftward along the
  bottom.
\item $a+2b+1+\gamma$ is an endpoint of colour $\gamma$.
\end{itemize}
A \emph{state} is the sorted list of the pieces, each given by its unordered pair of end codes,
together with a two-bit mask of the colours whose path is complete. The start state has no piece
and mask $0$. In the step for cell $(r,c)$, \emph{down} is the code $c$ if $r+1<a$ and exit slot
$a+b-1-c$ otherwise, and \emph{right} is the code $b$ if $c+1<b$ and exit slot $r$ otherwise.

% lean: ZZN/Statement.lean:573 stepCell
% lean: ZZN/Statement.lean:566 closeP
\begin{algorithm}[H]
\caption{Step of cell $(r,c)$ for one state $S$}
\begin{algorithmic}[1]
\State $u \gets$ the piece of $S$ with end code $c$, if $r>0$
\State $\ell \gets$ the piece of $S$ with end code $b$, if $c>0$
\State $\mathit{deg} \gets 1$ if $(r,c)$ is an endpoint, else $2$
\State $\mathit{in} \gets$ the number of $u,\ell$ present
\If{$\mathit{deg}<\mathit{in}$ \textbf{or} $u$ and $\ell$ are the same piece}
  \State \Return no successor
\EndIf
\State $E \gets$ the other end of $u$, the other end of $\ell$, and the endpoint code of $(r,c)$,
  those present
\For{each list $O$ of new ends: $[\,]$ if $\mathit{deg}=\mathit{in}$, $[\textit{down}]$ and
  $[\textit{right}]$ if $\mathit{deg}=\mathit{in}+1$, $[\textit{down},\textit{right}]$ otherwise}
  \If{$E$ followed by $O$ is exactly two ends $e_1,e_2$}
    \If{$e_1$ and $e_2$ are both endpoints}
      \If{they have the same colour and that colour is not complete}
        \State emit $S$ without $u,\ell$, with that colour marked complete
      \EndIf
    \Else
      \State emit $S$ without $u,\ell$, with the piece $\{e_1,e_2\}$ added
    \EndIf
  \EndIf
\EndFor
\end{algorithmic}
\end{algorithm}

% lean: ZZN/Statement.lean:619 normS
% lean: ZZN/Statement.lean:624 stage
% lean: ZZN/Statement.lean:657 stageF
% lean: ZZN/Statement.lean:661 dpAcceptF
% lean: ZZN/WinCheck.lean:122 inside
A \emph{stage} processes one cell for every state, sorts the successors by mask and pieces with a
merge sort, and removes duplicates. On a whole grid no piece can end in an exit slot, so the stage
also drops every state with an exit end. The instance is accepted when, after the last cell, the
state with no piece and mask $3$ is present. On a window the exit ends are kept, and the states
after the last cell are the boundary states of Section~3.2.

% lean: ZZN/Constr/ThinWit2.lean:26 stageW2
% lean: ZZN/Constr/ThinWit2.lean:38 thinSolveW2
% lean: ZZN/Timed/ThinRope.lean:149 thinSolveWRT
The version with paths carries, for every state, the cells of each open piece and the completed
paths, keeps one entry per state when it removes duplicates, and returns the two paths of the
accepting entry, transposed back if the instance was transposed. The counted version stores the
cells of a piece as a tree with constant-time append and reversal.

% lean: ZZN/Timed/ThinWord.lean:379 wcnt
% lean: ZZN/Timed/ThinRopeCost.lean:31 Nw
% lean: ZZN/Timed/ThinWordCount.lean:248 card_fold
% lean: paper/tools/KwrEval.lean (#eval at b = 2, 3); paper/tools/kwr.py (b = 10, 16)
\paragraph{The state bound $N(b)$.} Let $N(b)$ be the number of words of $b+1$ letters, one letter
per plug code $0,\ldots,b$, over the letters empty, open bracket, close bracket, end of colour $0$
and end of colour $1$, in which the brackets are balanced and each colour has at most two end
letters, each word counted twice for every colour with no end letter. Lean proves that a list of
distinct states reached after any number of cells has at most $N(b)$ entries, by mapping each
reached state injectively to such a word and its mask. Lean evaluates $N(2)=64$ and $N(3)=202$.
The transcription \lean{paper/tools/kwr.py} gives $N(10)=966\,244$ and
$N(16)=1\,381\,920\,144$.

\section{Proof of the separation fact}\label{sec:sep}

This section proves the separation fact used by the proofs of T1 and T2 in Section~3.2.

% lean: ZZN/Planar.lean:25 AdjZ
% lean: ZZN/Planar.lean:29 steps
% lean: ZZN/Planar.lean:31 Closed
% lean: ZZN/Planar.lean:42 ray
% lean: ZZN/Planar.lean:48 par
% lean: ZZN/Planar.lean:137 parP
Points are pairs $(i,j)$ of integers, $i$ the row and $j$ the column, and a unit step joins points
at distance $1$. A closed walk $\Gamma$ is a list of points in which consecutive points, and the
last and the first, are at unit distance. The unit square $F(i,j)$ has the corners $(i,j)$ and
$(i+1,j+1)$. Its \emph{crossing parity} $\pi(i,j)$ is the parity of the number of steps of $\Gamma$
between $(i',j)$ and $(i',j+1)$ with $i'\le i$, the steps that the upward ray from the centre of
$F(i,j)$ crosses. The crossing parity of a point $p=(i,j)$ is $\pi(i,j)$.

% lean: ZZN/Planar.lean:184 parP_chain
\begin{lemma}[separation]\label{lem:sep}
All points of a chain of points at unit distance, none of them on a closed walk $\Gamma$, have the
same crossing parity.
\end{lemma}
\begin{proof}
% lean: ZZN/Planar.lean:61 par_down
Write $e(\cdot)$ for the number of steps of $\Gamma$ on an edge, modulo $2$. The ray from $F(i+1,j)$
crosses the steps that the ray from $F(i,j)$ crosses and the steps on their common edge, so
$\pi(i+1,j)=\pi(i,j)+e(\text{common edge})$.

% lean: ZZN/Planar.lean:71 step_identity
% lean: ZZN/Planar.lean:96 par_right
For $F(i,j)$ and $F(i,j+1)$, with common edge from $(i,j+1)$ to $(i+1,j+1)$, let $V$ be the set
of points $(i',j+1)$ with $i'\le i$. A case check over the unit steps $(u,v)$ shows that the number
of the following that hold is even. The step crosses the ray of $F(i,j)$. It crosses the ray of
$F(i,j+1)$. It lies on the common edge. $u\in V$. $v\in V$. Summed over the steps of $\Gamma$, the
count of steps with $u\in V$ equals the count with $v\in V$, since every point of $\Gamma$ is the
first point of one step and the second point of another. So
$\pi(i,j+1)=\pi(i,j)+e(\text{common edge})$.

% lean: ZZN/Planar.lean:140 faces_around
% lean: ZZN/Planar.lean:162 parP_adj
If a point is not on $\Gamma$, no step of $\Gamma$ lies on the four edges that meet at it, so the
four squares around it have the same crossing parity. Two points at unit distance share two of
their four squares, so two such points that are not on $\Gamma$ have the same crossing parity.
Induction along the chain completes the proof.
\end{proof}

\section{Why one window certificate covers every frame, colouring and grid size}\label{sec:transfer}

This section proves the reduction behind the window certificates of Section~3.2, from a match of
B, C or D in any frame, colouring and wide grid to a fixed list of certificates.

% lean: ZZN/WinBCD.lean:40 CertFacts
% lean: ZZN/WinBCD.lean:22 winC
% lean: ZZN/WinBCD.lean:26 tmC
% lean: ZZN/WinBCD.lean:30 tmB
% lean: ZZN/WinBCD.lean:34 fsB
The certificates are stated in the identity frame. For C and D there is one certificate per
configuration of window side $6$, for the colours as listed and the parity $wh \bmod 2$ of its list.
For B there is one certificate per triple $(p,q,b)$, with colour $0$ at $p,q$ and colour $1$ at $b$,
the parity of its list, and each of three classes of the fourth endpoint $F$. The classes are
$F=(0,6)$, $F=(6,0)$ and every other cell. Each B certificate fixes the sign of $F$ to
\[
\sigma_F \;=\; 2\cdot(wh \bmod 2) - \sigma(p) - \sigma(q) - \sigma(b),
\]
and it is required only when this value is $\pm 1$, and for the classes $(0,6)$ and $(6,0)$ only
when it is $+1$. The six C configurations with a cell at coordinate $6$ or $7$ have no certificate.
Section~3.2 covers them by R.

% lean: ZZN/TheoremB.lean:39 frames_sound
% lean: ZZN/TheoremB.lean:34 frameMap_dims
% lean: ZZN/WinBCD.lean:50 solvable_swap
% lean: ZZN/WinBCD.lean:98 cd_sound
% lean: ZZN/WinBCD.lean:227 b_core
% lean: ZZN/WinBCD.lean:359 b_cert
% lean: ZZN/WinBCD.lean:388 bSound_of
% lean: ZZN/C8R.lean:363 cSound_of
\begin{lemma}[certificate transfer]\label{lem:transfer}
If B, C or D fires on a solvable wide instance, then some window certificate of the list above
applies to a solvable wide instance and gives a contradiction, except for the six C configurations
with a cell at coordinate $6$ or $7$.
\end{lemma}
\begin{proof}
\emph{Frames.} If the entry fires in frame $f$, it fires in the identity frame of the instance $J$
obtained by applying the reflections and the transposition of $f$. The instance $J$ is well-formed
and solvable, both its sides are at least $11$, and its area is that of the given instance.

\emph{Colours.} C and D match a configuration with either assignment of the two colours to its two
pairs, and B either colour at the pair $p,q$. Exchanging the colours keeps an instance well-formed
and solvable, so the colours can be taken as in the certificate.

\emph{The outside endpoint.} For B the completion test reads $F$ only through its colour, its sign
and the exit slot whose outside cell it is, if any. On a grid with both sides at least $11$, the
only outside cells of the $6\times 6$ window that lie on the border are $(0,6)$, outside slot $0$,
and $(6,0)$, outside slot $11$. Lemma~P on the solvable instance gives
$\sigma(F)=\sigma_F$ as above. Every cell has sign $\pm 1$, and $(0,6)$ and $(6,0)$ have sign $+1$.

% lean: ZZN/WinOut.lean:21 WinSetup
% lean: ZZN/WinSound.lean:24 window_sound
\emph{The window.} For C and D all four endpoints lie in the window, and for B three do and $F$
lies on the border outside it. The window side is below both sides of the grid. These are the
hypotheses of the window lemma of Section~3.2, which turns the certificate into a contradiction
with solvability.
\end{proof}

\section{Certificate formats and the programs that found them}\label{sec:certs}

This section gives the data and the checks of the window certificates of Section~3.2 and the grid
certificates of Section~3.4, and the programs that produced them.

\subsection{Window certificates}

% lean: ZZN/WinCheck.lean:178 FSpec
% lean: ZZN/WinCheck.lean:220 certify (defined for any sides a, b)
% lean: ZZN/WinBCD.lean:40 CertFacts (c_even only for winC AB = 6, b_even and b_odd with sides 6 6)
% lean: ZZN/PassX.lean:303 corner4_odd_small (every D configuration lies in the window of side 6)
A window certificate names the endpoints inside the window with their colours, the parity
$wh\bmod 2$, and for B the colour, sign and exit slot (or none) of the outside endpoint $F$. It
states that every boundary state of the window DP of Section~\ref{sec:plugdp} fails the completion
test of Section~3.2. The check is defined for a window of any side, and every certificate of the
list of Section~\ref{sec:transfer} has side $6$. The completion test is computed as follows.

% lean: ZZN/WinCheck.lean:181 usedOf
% lean: ZZN/WinCheck.lean:187 linkOf
% lean: ZZN/WinCheck.lean:200 fxOf
% lean: ZZN/WinCheck.lean:206 parOf
% lean: ZZN/WinCheck.lean:133 ncm
% lean: ZZN/WinCheck.lean:151 walk
% lean: ZZN/WinCheck.lean:162 visC
% lean: ZZN/WinCheck.lean:170 structOK
% lean: ZZN/WinCheck.lean:210 outsideOK
\begin{itemize}
\item The outside ends are numbered $0,\ldots,m-1$, first the used exit slots in increasing slot
  order and then $F$. Each outside end has an inside link. For a used slot it is the outside end
  at the other end of its piece, or the colour of the endpoint there. For $F$ it is its colour.
  The test fails if some piece at a used slot ends at a plug.
\item The test requires $m$ even and nonzero, and $\sigma$ summed over the cells just outside the
  used slots, plus $\sigma(F)$, equal to $2\cdot(wh\bmod 2)$.
\item The non-crossing perfect matchings of $0,\ldots,m-1$ are enumerated recursively. End $0$ is
  matched with an end $j$ such that an even number of ends lies between them, and the ends between
  and the ends after $j$ are matched separately.
\item A matching is accepted when the following hold. If $F$ is the outside cell of a used slot,
  that slot is matched with $F$. A colour completed inside the window has no outside end linked to
  its colour. Every other colour has exactly two such ends, and the walk from the first, alternately
  along the matching and along the inside links, ends at the second. Every outside end lies on one
  of these walks.
\item A boundary state satisfies the test when some matching is accepted.
\end{itemize}

% lean: ZZN/WinK/Engine.lean:1 (module docstring: checkpointed window certificates)
% lean: ZZN/WinK/Shared.lean:1 (module docstring: shared late stages)
% lean: WinKU/Facts.lean:167 certFactsK
The kernel checks each certificate through checkpoints. The set of states after every cell is
given as data, and the kernel checks that every successor of a state in one checkpoint lies in the
next and that every state of the last checkpoint fails the completion test. A stage sees the
endpoints only through the endpoint at its own cell. So for the cells after the last cell that
holds an endpoint in any certificate, all certificates share one checkpoint per cell, the union of
their state sets. The program \lean{certgen/tools/WinGenU.lean} computes the checkpoints for the
cases of the list \lean{certgen/tools/wincases.txt} that have side $6$ and the colours of the
certificates. The certificates involve no search.

\subsection{Grid certificates}

% lean: ZZN/Moves.lean:173 Sym
% lean: ZZN/BoxK/Box.lean:25 rowsOf
% lean: ZZN/BoxK/Box.lean:26 colsOf
% lean: ZZN/BoxK/Slice.lean:191 keyR
% lean: ZZN/BoxK/Rows.lean:25 G16
% lean: ZZN/BoxK/Box.lean:54 Gbox
The \emph{symmetries} are generated by reflecting the $x$-axis, reflecting the $y$-axis,
transposing, exchanging $s_0$ with $t_0$, exchanging $s_1$ with $t_1$, and exchanging the colours.
The \emph{row tuple} of an instance lists the $x$-coordinates of $s_0,t_0,s_1,t_1$, and the
\emph{column tuple} their $y$-coordinates. The \emph{rank} of a tuple $(a,b,c,d)$ with entries
below $32$ is $((32a+b)32+c)32+d$, which orders tuples lexicographically.

% lean: ZZN/BoxK/Box.lean:142 BoxCert
% lean: ZZN/BoxK/Rows.lean:60 rowsOK
% lean: ZZN/FiniteB.lean:36 axBB
% lean: ZZN/BoxK/Box.lean:130 SliceGood
% lean: ZZN/BoxK/Slice.lean:224 stabL
% lean: ZZN/BoxK/Slice.lean:247 trV
% lean: ZZN/BoxK/Slice.lean:70 distinctV
% lean: ZZN/BoxK/Slice.lean:98 parityV
% lean: ZZN/BoxK/Slice.lean:355 sliceOK
% lean: ZZN/BoxK/SliceC.lean:138 sliceCOK
% lean: ZZN/BoxK/SliceC.lean:144 CData
% lean: ZZN/BoxK/FireW.lean:27 fireW
% lean: ZZN/BoxK/FireW.lean:57 failW
% lean: ZZN/Finite.lean:411 passImgB
% lean: ZZN/BoxK/Slice.lean:335 solsM
% (sliceCOK also accepts a chain from a smaller slice and a mask M of lanes solved there; in all
%   16 477 slices of proof/BoxKGen w0 = w and M = 0, so the text does not describe them.)
A grid certificate for the size $w\times h$, $11\le w\le h\le 16$, consists of a list of
\emph{row classes} and data for each row class. The list contains every row tuple with entries
below $w$ on which no pair deletion applies on the $x$-axis and whose rank is least among its
images under the $16$ symmetries generated by reflecting the $x$-axis and the three exchanges. The
\emph{covered column tuples} of a row class are those with entries below $h$ on which no pair
deletion applies on the $y$-axis, whose four cells are distinct, that satisfy P, whose rank is
least among their images under the symmetries without transposition that fix the row tuple, and,
when $w=h$, whose rank is at least that of the row tuple. The data of a row class are the
following.
\begin{itemize}
% lean: ZZN/Cut/FastW.lean:231 coverFWD (pairs count only without transposition or y-reflection,
%   and with small offsets, tsmall)
\item Pairs of a template and a symmetry. The check counts a pair for a covered instance only
  when the symmetry is one of the $16$ symmetries above and the template checks on the image of
  the instance under it.
\item Failing records, each a column tuple with a witness. The witness names, for the instance and
  for its images under reflecting the $x$-axis, the $y$-axis and both, one entry that fires there.
\item Explicit pairs of paths, each for one column tuple. They are checked to be a solution.
\end{itemize}
The check states that every covered column tuple is solved by a pair, recorded as failing, or
solved by explicit paths. It is evaluated for all $h^4$ column tuples of a row class at once, on
bit lanes of large natural numbers. A checked row class gives that every covered instance is
solvable or fails the catalogue together with its three reflections.

% lean: ZZN/PassSym.lean:853 passes_rev0, :858 passes_rev1, :863 passes_swapColors
% lean: ZZN/PassSym.lean:671 frameFires_perm, :680 frameFires_relab, :709 t2_swap01, :738 t2_pairs
% lean: ZZN/Orbit.lean:249 passes_lab
% lean: ZZN/Orbit.lean:274 passes_transpose_iff
% lean: ZZN/Orbit.lean:282 passes_refl_of_image
On a wide instance, exchanging $s_0$ with $t_0$, exchanging $s_1$ with $t_1$ and exchanging the
colours each keep the answer of every entry, which Lean proves entry by entry. Transposition
keeps passing by Section~\ref{sec:rtrans}. So if some image of a wide instance under a symmetry passes, then the
instance or one of its three reflections passes.

\subsection{The search programs}

% source: certgen/src/chain3.c:1-5, 127-135 (near), 204-240 (solve), 292-297 (iterative deepening)
% source: certgen/src/tmpl.c:1-10, 27-41 (enc: nearest anchor), 141-145 (anchors)
% source: certgen/src/tx.c:1-5; certgen/src/exh2.c:1-6; certgen/src/peel.c:1-6
% source: certgen/src/slicehyb.c:1-9, 348-351 (SOLS); certgen/regen.sh:6-8, 22-31;
%   certgen/tools/gen_cbox.py; certgen/tools/WinGenU.lean:16 (side A, names ending in c0 or A0)
The programs in \lean{certgen/src} found the templates and the grid certificates. They are outside
Lean, and the kernel checks their output.
\begin{itemize}
\item \lean{chain3.c} searches by iterative deepening for a cut tree of least depth. A rectangle
  holds at most three terminal pairs, every rectangle passes the parity test, and cut and crossing
  positions lie within one line of a terminal of the rectangle or of its boundary. Leaves use a
  transcription of the IPS criterion.
\item \lean{tmpl.c} turns a found tree into a template by writing each position as an offset from
  the nearest of its six anchors (the four endpoint coordinates on that axis, $0$, and the side
  length), and tries the known templates before searching. \lean{tx.c} folds the reflections of a
  symmetry into a template. \lean{exh2.c} and \lean{peel.c} test how far a template set covers the
  instances of a grid size.
\item \lean{slicehyb.c} builds the certificate of one row class. It collects the
  (template, symmetry) pairs that solve covered instances, orders them by a greedy cover, records
  the column tuples whose instance and three reflections all fail the catalogue, and solves each
  remaining instance by the tree search, whose tree becomes a new template, or, in a separate
  mode, by explicit paths from the C reference solver.
\item \lean{certgen/regen.sh} writes the Lean modules of \lean{proof/BoxKGen} from the committed
  certificates in \lean{certgen/certs}, without searching again, and those of \lean{proof/WinKU}
  from the case list through \lean{WinGenU.lean}.
\end{itemize}

\section{Pair deletions keep each entry passing}\label{sec:keep}

This section gives the cases behind the lemma of Section~3.3 that a pair deletion turns a passing
wide instance into a passing wide instance.

% lean: ZZN/Skeleton.lean:50 RedX (clauses 1, 2: edge strip. Clause 3: wide-gap compression.
%   Clause 4: the distant-endpoint compressions with 8 <= d and d + 10 <= w)
% lean: ZZN/PassB.lean:29 RedB (distant-endpoint compression), ZZN/PassX.lean:689 CompCore
% lean: ZZN/PassB.lean:181 passX_all
% lean: ZZN/PassB.lean:118 passB
% lean: ZZN/PassXAssemble.lean:20 comp_canonical
% lean: ZZN/PassXR.lean:25 deleteAt_window
The proof shows that an entry firing on the smaller instance $I'$ makes an entry fire on the
larger instance $I$, for each of the three kinds of pair deletion, the edge strip, the wide-gap
compression and the distant-endpoint compression. A wide-gap compression of a run
$x=\ell,\ldots,x=\ell+9$ of free lines gives the same $I'$ as deleting its two middle lines
$x=\ell+4$ and $x=\ell+5$, for which $5\le\ell+4$ and $\ell+4+7\le w$. The proof takes it at these
lines. So every compression below has $5\le d$ and $d+7\le w$. Distances between lines are
differences of coordinates, and the frame coordinates of a cell are its distances to the two
border lines at the frame's corner.

% lean: ZZN/PassX.lean:40 parityOk_deleteAt
\paragraph{P.} Deleting two lines changes neither $wh\bmod 2$ nor the sign of any endpoint, since
moved endpoints move by $2$.

% lean: ZZN/PassX.lean:1305 t1_comp
% lean: ZZN/PassX.lean:895 t2_comp
% lean: ZZN/PassX.lean:1117 l1_comp
% lean: ZZN/PassX.lean:980 l6_comp
% lean: ZZN/PassStrip.lean:29 NearP
\paragraph{T1, T2, L1, L6.} These depend on the clockwise order of the endpoints on the border, on
which endpoints are adjacent, and on which cells next to a corner hold endpoints. An edge strip
moves every endpoint by the same amount, $0$ or $2$, and leaves at least $2$ free lines at the
stripped edge. A compression deletes lines between endpoint lines, leaves a free line between the
endpoint lines on its two sides, and, with $5\le d$ and $d+7\le w$, keeps every distance below $5$
between an endpoint and a border line. In both cases these features agree on $I$ and $I'$.

% lean: ZZN/PassStrip.lean:82 boundary_lists_local
% lean: ZZN/PassStrip.lean:89 corner4_lists_local
% lean: ZZN/PassStrip.lean:204 edgeClosure_flip
% lean: ZZN/PassStrip.lean:1280 effAlt3_lstrip
% lean: ZZN/PassStrip.lean:608 effAlt3_rstrip
\paragraph{Edge strip, frame entries and R.} In $I'$ every endpoint is at least $2$ lines from the
stripped edge. In the four frames whose corner lies on that edge, the first two frame rows are
then empty, or the first two frame columns for a transposed frame. L2 to L5 and every listed
configuration of B, C and D need an endpoint in the first two rows and one in the first two
columns, and E1 needs endpoints in the first two rows. So in these four frames only E1 in a
transposed frame can fire, and its reflection along the stripped edge is E1 with the colours
exchanged in a frame whose corner is not on that edge. In the other four frames the endpoints keep
their frame coordinates in $I'$. R gives the same answer on $I$ and $I'$, since in $I$ every
endpoint is at least $4$ lines from the stripped edge.

% lean: ZZN/PassX.lean:155 FarK
% lean: ZZN/PassX.lean:68 GoodInK, :379 WinRelK (border in I' implies border in I)
% lean: ZZN/PassX.lean:774 winRel_comp
% lean: ZZN/PassX.lean:620 frames_any_of
% lean: ZZN/PassX.lean:848 e1_comp
% lean: ZZN/PassXAssemble.lean:40 passX_comp
\paragraph{Wide-gap compression, frame entries.} Every endpoint keeps its distance to each border
line or has distance at least $8$ to it in both $I$ and $I'$. So the frame cells with both
coordinates below $8$ hold the same endpoints in both instances, and an endpoint that lies on the
border, or outside the $6\times 6$ window of a frame, in $I'$ does so in $I$. L2 to L5, B, C and D
read only these facts and $wh\bmod 2$. E1 along the deleted axis would need four endpoints on four
consecutive lines. Four such lines on one side of the deleted pair are consecutive in $I$ too, and
four lines on both sides are impossible in $I'$, where a free line separates the two sides. The
same argument applies to a distant-endpoint compression with $8\le d$ and $d+10\le w$.

% lean: ZZN/PassX5.lean:180 frames_any_of5
% lean: ZZN/PassB.lean:77 bigHit_R
% lean: ZZN/PassB.lean:105 b5Hit_R
% lean: ZZN/C8R.lean:286 c8_firesF
% lean: ZZN/B5Fire.lean:78 b5_all
% lean: ZZN/PassX5.lean:99 WinRel5, :217 winRel5_core
% lean: ZZN/PassX5.lean:27 small6_le4; ZZN/PassX.lean:307 BigHit; ZZN/PassX5.lean:93 B5Hit
\paragraph{Distant-endpoint compression, frame entries.} The frame cells with both coordinates
below $5$ hold the same endpoints in $I$ and $I'$, and an endpoint on the border of $I'$ is on the
border of $I$. The frame coordinate of an endpoint along the deleted axis is, in $I$, its
coordinate in $I'$ or that plus $2$, so the test of B that $F$ lies on the border outside the
window carries over from $I'$ to $I$. Every listed configuration has all its cells below coordinate
$5$ except the six C configurations marked $*$, with a cell at coordinate $6$ or $7$, and the two
B triples marked $\dagger$, (01 05 11) and (05 11 10). These belong to the lists for even $wh$, and
a match of one of them on the wide instance $I'$ makes R fire on $I'$, by the paper's lemma on
configurations marked $*$ and $\dagger$. The last paragraph of this section outlines its proof.
Since R gives the same answer on $I$ and $I'$, R fires on $I$. E1 is excluded as for the wide-gap
compression.

% lean: ZZN/PassXR.lean:137 Av
% lean: ZZN/PassXR.lean:874 effAlt3_comp
\paragraph{R, compressions.} R gives the same answer on $I$ and $I'$ whenever the lines $x=d$ and
$x=d+1$ hold no endpoint, $5\le d$, $d+7\le w$ and $h\ge 4$. This covers both kinds of compression.
With $5\le d$ and $d+7\le w$ the deleted lines lie outside the four rows and columns at each
corner that the rules read and change, and the bypass walks of $I$ and $I'$ visit the effective
endpoints in the same order.

% lean: ZZN/B5R.lean:29 b5q
% lean: ZZN/B5R.lean:32 perimOut
% lean: ZZN/B5R.lean:68 perimOut_step
% lean: ZZN/B5R.lean:81 b5_stepR
% lean: ZZN/B5Fire.lean:71 b5_transpose
% lean: ZZN/B5Fire.lean:78 b5_all
% lean: ZZN/B5Base.lean:511 b5_base
% lean: ZZN/PassB.lean:105 b5Hit_R
% lean: ZZN/PassSym.lean:186 effAlt3_relab
% lean: ZZN/C8R.lean:160 effAlt3_redF, :267 c8_baseF, :286 c8_firesF
% lean: ZZN/PassStrip.lean:608 effAlt3_rstrip, :1280 effAlt3_lstrip
\paragraph{The configurations marked $*$ and $\dagger$.} Let one of the triples (01 05 11) and (05 11 10) match
in a frame of a wide instance with $wh$ even, with the fourth endpoint $F$ on the border outside
the $6\times 6$ window of the frame. Then R fires. Exchanging the colours keeps the answer of R,
so the two cells of the triple with the same colour are taken of colour $0$. The proof is an
induction on $w+h$. If the side
along the $x$-axis is at least $15$, let $c$ be the frame coordinate of $F$ on that axis, and
delete the two lines at frame coordinates $8$ and $9$ if $c$ is $6$ or $7$, and at frame
coordinates $6$ and $7$ otherwise. These two lines hold no endpoint, since the cells of the triple
have frame coordinates at most $5$ and $F$ is on neither by the choice of the pair, and they are lines $x=d$, $x=d+1$ with $5\le d$ and $d+7\le w$.
So R gives the same answer on the smaller instance, by the paragraph above. There the triple
matches in the same frame and $F$ is on the border outside the window. If only the side along the
$y$-axis is at least $15$, the same step is made on the transpose, on which R gives the same answer
(Section~\ref{sec:rtrans}). When both sides are at most $14$, the kernel checks that R fires, for
both triples, the $8$ frames, every size with sides $11$ to $14$ and even area, and every cell $F$
on the border outside the window. The proof for the six C configurations marked $*$ has the same
form. All four endpoints have frame coordinates below $8$. While a side is at least $13$, its two
lines with the greatest frame coordinates are deleted, which keeps the answer of R as for an edge
strip, every endpoint being at least $4$ lines from that edge. This brings each side to $11$ or
$12$, and the kernel checks the sizes $11\times 12$, $12\times 11$ and $12\times 12$ in the $8$
frames, the size $11\times 11$ having odd area.

\section{The forced-route test under transposition}\label{sec:rtrans}

This section proves the transposition part of the lemma of Section~3.3, that a wide instance
passes if and only if its transpose passes.

% lean: ZZN/PassTR.lean:1802 passT
% lean: ZZN/PassT.lean:306 fires3_transpose
% lean: ZZN/PassTR.lean:1755 effAlt3_transpose
\begin{lemma}[transposition]\label{lem:trans}
On a wide instance every entry gives the same answer on the instance and on its transpose.
\end{lemma}
\begin{proof}
% lean: ZZN/PassT.lean:306 fires3_transpose
Transposition keeps signs and the area, maps the corner frames to the corner frames, keeps
adjacency, and reverses the clockwise order on the border, which keeps alternation of four
endpoints. So every entry other than R gives the same answer.

% lean: ZZN/PassTR.lean:1057 frames_tp_eq
% lean: ZZN/PassTR.lean:1060 rewrite_order
For R, the transpose visits the frames in a different order. At each corner the transposed frame
comes first, and the corners at $(0,h-1)$ and $(w-1,0)$ exchange places. Steps at different
corners commute, since a rule reads and changes only the $4\times 4$ square at its corner and
these squares are disjoint when both sides are at least $8$.

% lean: ZZN/PassTR.lean:337 rewritesAt_both
% lean: ZZN/Statement.lean:260 rewritesAt
At one corner, if the rules of both frames match, they are the same rule, R2 or R3, since R1 in one
frame needs an endpoint at a cell that every rule of the other frame needs empty. R3 makes the same
two moves and settles the same cells in both frames. R2 settles the same four cells in both frames
and moves the two endpoints of the colour at $(0,0)$ and $(1,1)$ to the same two cells, with the
targets of the two endpoints exchanged. Every rule settles the corner cell, so the frame visited
second at a corner applies no rule once the first has applied one. So the settled cells and the
set of effective endpoints with their colours do not depend on the order.

% lean: ZZN/PassTR.lean:1281 walk3_swap
% lean: ZZN/PassTR.lean:1578 remShape_rewriteAll
On a wide instance the settled cells at each corner form one of a fixed list of shapes. For these
shapes the bypass walk of the transpose is the transposed bypass walk $T\,Y$ written as
$T^{\mathrm r}\,Y^{\mathrm r}$, each of two consecutive parts read in reverse, which is the reverse
of the walk $Y\,T$ started at another cell. With two endpoints of each colour, whether the colours
alternate along a closed walk depends neither on the starting cell nor on the direction, and the
number of times each effective endpoint occurs on the walk is unchanged. So R gives the same
answer.
\end{proof}

\section{Block selection, block lifts and table lookup}\label{sec:blocks}

This section gives the procedures that Section~3.5 names, SelectBlock, EdgeSplice, MiddleSplice
and TableLookup. The blocks and lifts are stated on the $x$-axis. On the $y$-axis they are applied
to the transpose.

% lean: ZZN/Constr/Round.lean:271 xmin
% lean: ZZN/Constr/Round.lean:275 leftK
% lean: ZZN/Constr/Round.lean:277 rightK
% lean: ZZN/Constr/Round.lean:279 innerK
% lean: ZZN/Constr/Round.lean:282 gapOK
% lean: ZZN/Constr/Round.lean:286 innerPair
% lean: ZZN/Constr/Round.lean:290 blkX
% lean: ZZN/Constr/Round.lean:50 Inst.dropIn
% lean: ZZN/Finite.lean:131 xsOf
\begin{algorithm}[H]
\caption{SelectBlock$(I)$}\label{alg:select}
\begin{algorithmic}[1]
\Require a wide instance $I$ with $x$-axis length $n$
\Ensure a block of $2k$ free lines to delete, or none
\Statex \hspace{\algorithmicindent}(a subtraction with a negative result is taken as $0$)
\State $x_{\min}, x_{\max} \gets$ the least and greatest endpoint $x$-coordinate
\State $m \gets \lfloor (n-11)/2\rfloor$
\State $k \gets \min(\lfloor (x_{\min}-2)/2\rfloor, m)$
\If{$k\ge 1$} \Return the top block, lines $0$ to $2k-1$ \EndIf
\State $k \gets \min(\lfloor (n-3-x_{\max})/2\rfloor, m)$
\If{$k\ge 1$} \Return the bottom block, the last $2k$ lines \EndIf
\For{$a$ in the $x$-coordinates of $s_0,t_0,s_1,t_1$, in this order}
  \For{$b$ in the $x$-coordinates of $s_0,t_0,s_1,t_1$, in this order}
    \State $k \gets \min(\lfloor (b-a-9)/2\rfloor, m)$
    \If{$a<b$ \textbf{and} no endpoint has $a<x<b$ \textbf{and} $k\ge 1$}
      \State \Return the middle block, lines $a+2$ to $a+2k+1$
    \EndIf
  \EndFor
\EndFor
\State \Return none
\end{algorithmic}
\end{algorithm}

% lean: ZZN/Constr/Round.lean:322 leftK_spec
% lean: ZZN/Constr/Round.lean:329 rightK_spec
% lean: ZZN/Constr/Round.lean:336 innerK_spec
% lean: ZZN/Constr/Round.lean:344 blkX_ok
A top or bottom block leaves at least $2$ free lines between it and the nearest endpoint line, a
middle block leaves at least $8$ free lines between $a$ and $b$, and every block leaves the side at
least $11$. So deleting a block is a sequence of pair deletions, edge strips or wide-gap
compressions.

% lean: ZZN/Constr/EdgeBlock.lean:33 colSnake
A \emph{snake} over a list $A$ of columns on the new lines $x_0,\ldots,x_0+2k-1$ runs along $A$ on
line $x_0$, steps to line $x_0+1$, runs back along $A$, and so on through all $2k$ lines. It
starts at $(x_0,A_1)$ and ends at $(x_0+2k-1,A_1)$, where $A_1$ is the first column of $A$.

% lean: ZZN/Constr/EdgeBlock.lean:181 edgeBlock
% lean: ZZN/Constr/EdgeBlock.lean:300 findV
% lean: ZZN/Constr/EdgeBlock.lean:339 blkE
% lean: ZZN/Constr/EdgeBlock.lean:374 edgeLift
% lean: ZZN/Constr/EdgeBlock.lean:383 exists_vert
% lean: ZZN/Constr/EdgeBlock.lean:413 edgeLift_sol
\begin{algorithm}[H]
\caption{EdgeSplice$(p,q,k)$, a bottom block of $2k$ lines}\label{alg:edge}
\begin{algorithmic}[1]
\Require a solution $p,q$ of a $w\times h$ instance with $h\ge 5$
\Ensure a solution of the instance with $2k$ lines added after the line $x=w-1$
\State find the first step of $p$, else of $q$, between cells $(w-1,y)$ and $(w-1,y+1)$, in
  either direction
\State $S \gets$ a snake over the columns $y, y-1, \ldots, 0$ on the new lines, followed by a snake
  over the columns $y+1,\ldots,h-1$ walked backward, from $(w,y)$ to $(w,y+1)$
\State insert $S$ into that step, reversed if the step runs from $(w-1,y+1)$ to $(w-1,y)$
\State \Return the two paths
\end{algorithmic}
\end{algorithm}

% lean: ZZN/Constr/EdgeBlock.lean:383 exists_vert
The step exists because the last line has at least $5$ cells, so one of them, say $(w-1,y_0)$, is
not an endpoint. That cell has two path neighbours, at most one of them on the line $x=w-2$, so its
path steps along the last line. A top block is put back by reflecting the $x$-axis, applying
EdgeSplice, and reflecting back.
% lean: ZZN/Constr/EdgeBlock.lean:572 edgeLiftL
% lean: ZZN/Constr/EdgeBlock.lean:577 edgeLiftL_sol

% lean: ZZN/Statement.lean:114 Cross
% lean: ZZN/Statement.lean:425 SeedAt
% lean: ZZN/Constr/Round.lean:202 seedW
% lean: ZZN/Constr/Round.lean:208 cvW
% lean: ZZN/Constr/WideBand.lean:278 attW
% lean: ZZN/Constr/WideBand.lean:281 grpW
% lean: ZZN/Constr/WideBand.lean:285 snW
% lean: ZZN/Constr/WideBand.lean:288 flipW
% lean: ZZN/Constr/WideBand.lean:292 GW
% lean: ZZN/Constr/WideBand.lean:300 wideExt
% lean: ZZN/Constr/Round.lean:212 innerLift
MiddleSplice uses the crossings and seeds of the band-extension lemma of Section~3.3. Column $y$ is
a \emph{crossing} when a path steps between $(r,y)$ and $(r+1,y)$. Position $R$, $0\le R\le h$, lies
between columns $R-1$ and $R$. It is a \emph{gap seed} when column $R-1$ is a crossing or $R=0$,
and column $R$ is a crossing or $R=h$. It is a \emph{step seed} when $1\le R<h$, neither column
$R-1$ nor column $R$ is a crossing, and a path steps between $(r,R-1)$ and $(r,R)$.

\begin{algorithm}[H]
\caption{MiddleSplice$(p,q,r,k)$, a middle block of $2k$ lines after the free line $x=r$}\label{alg:mid}
\begin{algorithmic}[1]
\Require a solution $p,q$ of an instance in which the line $x=r$ is free, $r<w$
\Ensure a solution of the instance with $2k$ lines inserted after $x=r$
\State $R \gets$ the least position that is a gap seed or a step seed
\State $\mathit{flip} \gets$ true if $R$ is not a gap seed
\State the \emph{source} columns are the crossings, and $R-1$ and $R$ if $\mathit{flip}$
\State the \emph{group} of a source column $a$ is $a$ followed by the columns $a-1,a-2,\ldots$ if
  $a<R$, or $a+1,a+2,\ldots$ if $a\ge R$, stopping before the next source column or the edge of
  the grid
\State move the cells with $x>r$ down by $2k$
\For{each step of $p$ and $q$ between $(r,a)$ and $(r+2k+1,a)$}
  \State insert a snake over the group of $a$ on the new lines $r+1,\ldots,r+2k$, reversed when
    the step runs from $(r+2k+1,a)$ to $(r,a)$
\EndFor
\If{$\mathit{flip}$}
  \State insert into the step between $(r,R-1)$ and $(r,R)$ a snake over the group of $R-1$,
    followed by a snake over the group of $R$ walked backward, all reversed when the step runs from
    $(r,R)$ to $(r,R-1)$
\EndIf
\State \Return the two paths
\end{algorithmic}
\end{algorithm}

% lean: ZZN/Constr/WideBand.lean:856 wideExt_sol
% lean: ZZN/Constr/Round.lean:246 innerLift_sol
% lean: ZZN/BandExtension.lean:225 seed_or_flip
Every group extends away from the seed. Every column that is not a source column has a source
column on its seed side, and it belongs to the group of the nearest one only. So the snakes fill
every new cell once, and the result is a solution. A seed exists by the counting argument of the
band-extension lemma.

% lean: ZZN/Look/Leaf.lean:46 entryP
% lean: ZZN/Look/Leaf.lean:96 lookK
% lean: ZZN/Look/Leaf.lean:102 argminS
% lean: ZZN/Look/Leaf.lean:136 lookLe
% lean: ZZN/Look/Leaf.lean:141 lookJ
% lean: ZZN/BoxK/Box.lean:54 Gbox
% lean: ZZN/BoxK/Box.lean:94 key2
% lean: ZZN/Look/Sym.lean:137 unSym
% lean: ZZN/Look/Build.lean:428 tsol
% lean: ZZN/Look/Table.lean:64 tabAll
The table holds, for each grid size $11\le w\le h\le 16$, the row classes of its grid certificate
(Section~\ref{sec:certs}) with their pairs and explicit paths.

\begin{algorithm}[H]
\caption{TableLookup$(J)$}\label{alg:lookup}
\begin{algorithmic}[1]
\Require an instance $J$ with both sides between $11$ and $16$
\Ensure a solution of $J$, or none
\If{$w>h$} \State \Return the transpose of TableLookup$(J^{\mathsf T})$ \EndIf
\State $G\gets$ all symmetries if $w=h$, else the symmetries without transposition
\State $\rho\gets$ the first $\rho\in G$ that minimizes the rank of the row tuple of $\rho(J)$ and
  then the rank of its column tuple
\State $\hat J\gets\rho(J)$
\State find the first table entry of size $w\times h$ whose row class is the row tuple of $\hat J$,
  and \Return none if there is none
\For{each pair $(T,\tau)$ of the entry, in order}
  \If{$T$ checks on $\tau(\hat J)$}
    \State build the paths of $T$ on $\tau(\hat J)$ \Comment{Section~\ref{sec:ips} at the leaves}
    \State \Return them mapped back by $\tau^{-1}$ and then $\rho^{-1}$
  \EndIf
\EndFor
\If{the entry holds explicit paths for the column tuple of $\hat J$ that solve $\hat J$}
  \State \Return them mapped back by $\rho^{-1}$
\EndIf
\State \Return none
\end{algorithmic}
\end{algorithm}

% lean: ZZN/Constr/Solve2.lean:214 solveRec2
% lean: ZZN/Statement.lean:519 checkE
% lean: ZZN/Constr/Round.lean:420 solveR
% lean: ZZN/Constr/Round.lean:528 solveZ3_sound
% lean: ZZN/Constr/Round.lean:246 innerLift_sol
% lean: ZZN/Constr/EdgeBlock.lean:413 edgeLift_sol
% lean: ZZN/Skeleton.lean:207 lift_redX
% lean: ZZN/PassB.lean:50 lift_redB
In Construct the final solver's answer is checked, and every lift, of a block or of a pair
deletion, is proved to map a solution of the smaller instance to a solution of the larger one. The
lifted solutions are not checked again.

\section{Building a Hamiltonian path in a rectangle}\label{sec:ips}

This section gives the recursion that builds the Hamiltonian paths at the leaves of the templates
used by TableLookup in Section~3.5.

% lean: ZZN/Look/IPS.lean:402 ipsP
% lean: ZZN/Look/IPS.lean:395 stepP
% lean: ZZN/Look/IPS.lean:386 normP
% lean: GridHam/Reduction.lean:267 NotPrime
% lean: GridHam/Reduction.lean:400 rChoice
% lean: GridHam/PrimeTable.lean:652 primeShapes
% lean: GridHam/PrimeTable.lean:659 primeLookup
% lean: ZZN/Look/IPS.lean:246 stripRightP
% lean: ZZN/Look/IPS.lean:330 stripLeftP
% lean: ZZN/Look/IPS.lean:202 colEdge
% lean: GridHam/StripSplit.lean:770 rightDetour
The recursion is on $w+h$, for an IPS-admissible pair $s,t$ of a $w\times h$ rectangle~\cite{ips82}.
The rectangle is transposed if needed so that $h\le w$, and with $s=(s_x,s_y)$, $t=(t_x,t_y)$ the
cells $s$ and $t$ are exchanged when $s_x>t_x$, the result then being reversed. The shapes with both
sides at most $3$ and the shapes $4\times 4$, $4\times 5$ and $5\times 4$ are answered from a stored
% lean: GridHam/PrimeTable.lean:668 prime_check
table of paths, which the kernel checks. Otherwise the first of three steps that applies is taken.
\begin{itemize}
\item If $s_x\ge 2$, build a path on the rectangle without the lines $x=0,1$, then cover those two
  lines by a detour from the first step of the path along the line $x=2$.
\item Else, if $t_x+2<w$, build a path on the rectangle without the last two lines, then cover
  them by a detour from the first step of the path along the line $x=w-3$.
\item Else, split after the line $x=1$. Build a path from $s$ to $(1,r)$ on the first two lines and
  a path from $(2,r)$ to $t$ on the rest, and join them. Here $r=0$ if $s_x+s_y$ is even, $r=3$ if
  $s=(0,1)$ and $h>2$, and $r=1$ otherwise.
\end{itemize}
A detour replaces a step between $(x,y)$ and $(x,y+1)$ on the line next to the two stripped lines
by a run through all cells of the stripped lines.

% lean: ZZN/Look/IPS.lean:474 ipsP_valid
% lean: GridHam/Main.lean:171 ips_characterization
Lean proves that the result is a Hamiltonian path from $s$ to $t$ whenever the pair is
IPS-admissible and the recursion depth is at least $w+h$. The proof follows the formalized proof
of the IPS characterization, which gives the step along the line next to the stripped lines and the
admissibility of the smaller pairs.

\section{The counted model and the step accounting}\label{sec:steps}

This section gives the step model, the exact per-cell cost $K(b)$ of the counted plug DP, and the
origin of each term of the step bounds of Section~3.5.

\subsection{The counted model}

% lean: ZZN/TimedStatement.lean:23 TimeM
% lean: ZZN/TimedStatement.lean:34 lmapT
% lean: ZZN/TimedStatement.lean:36 lfilterT
% lean: ZZN/TimedStatement.lean:44 lfilterMapT
% lean: ZZN/TimedStatement.lean:46 lflatMapT
% lean: ZZN/TimedStatement.lean:48 lfoldlT
% lean: ZZN/TimedStatement.lean:62 rangeT
% lean: ZZN/TimedStatement.lean:63 ldedupT
% lean: ZZN/TimedStatement.lean:64 lnodupT
% lean: ZZN/TimedStatement.lean:68 fmemT
A counted procedure returns a value with a number of steps. Returning a value costs $0$ steps, and
running two counted procedures in sequence adds their steps. The charged operations are the
following.
\begin{itemize}
\item Map, filter, filter-map, all, any, find and left fold over a list cost, for every element,
  the callee's steps plus $1$. Find charges every element of the list. A flat map also pays the
  length of each result. A filter over $\{0,\ldots,n-1\}$ as a finite set is charged the same way.
\item Length, membership, erase, append, reverse, sum, indexing, zip, zip with indices, find index,
  modify, last element and the test for a two-element infix cost the length of the list plus $1$.
  Append and zip charge the first list.
\item The range $0,\ldots,n-1$ costs $n+1$. Membership, erase and filter on a finite set cost its
  size plus $1$.
\item Removing duplicates and testing for repeated elements cost the square of the length plus
  $1$.
\end{itemize}
Arithmetic and comparisons on natural numbers of any size, and operations on Booleans, cells,
endpoint records, frames and options, are not charged, except where a definition charges a fixed
amount. The bounds count the charged steps.

% lean: ZZN/Timed/DecideR.lean:37 solvableBRT_iff
% lean: ZZN/Timed/RoundT.lean:597 solveZ3T_ret
% lean: ZZN/Look/RoundTimeL.lean:30 leafLT_ret
% lean: ZZN/Timed/ThinRope.lean:149 thinSolveWRT
The counted decision procedure returns true exactly when the instance is solvable, and the counted
construction returns the answer of Construct, with the table lookup or with the plug DP as final
solver.

\subsection{The per-cell cost of the plug DP}

% lean: ZZN/Timed/ThinRopeCost.lean:34 LgW
% lean: ZZN/Timed/ThinRopeCost.lean:37 KsR
% lean: ZZN/Timed/ThinRopeCost.lean:41 KwR
% lean: ZZN/Timed/ThinRopeCost.lean:393 thinSolveWRT_time
% lean: ZZN/Timed/ThinRopeCost.lean:352 solveDPRT_time
% lean: ZZN/Timed/ThinRopeCost.lean:381 solveDPRT_le
With $N=N(b)$ from Section~\ref{sec:plugdp} and $\lambda(b)=\lfloor\log_2(2N)\rfloor+1$ merge-sort
levels, the cost of one stage is
\[
K_s(b) = 4 + N\,(9(b+1)+31) + 4N\,(b+3) + 2N + 2 + 2(b+8)\,N\,(\lambda(b)+1),
\]
and the per-cell cost is $K(b)=K_s(b)+2N+30$. The counted plug DP takes at most
$(wh+1)\,K(\min(w,h))$ steps. Beyond the stages, the term $2N+30$ per cell, taken $wh+1$ times,
pays for the loop over the cells, the search of the final states for the accepting one, reading
out and orienting the two paths, and transposing them back.

% values: paper/tools/kwr.py (a transcription of wcnt, Nw, LgW, KsR, KwR); Lean #eval at b = 2, 3
%   in paper/tools/KwrEval.lean
The values of $K(10)$ and $K(16)$ are computed outside Lean, by \lean{paper/tools/kwr.py}, a
transcription of the definitions. It gives $K(10)=944\,986\,668$ and
$K(16)=2\,553\,788\,426\,148$. It agrees with Lean's evaluation $K(2)=16\,804$ and
$K(3)=63\,666$.

\subsection{Decision}

% lean: ZZN/Timed/CatalogueCost.lean:628 fires3T_time
% lean: ZZN/Timed/DecideR.lean:35 CdecR
% lean: ZZN/Timed/DecideR.lean:52 solvableBRT_time
The counted catalogue takes at most $300\,000\,(w+h+1)$ steps, and the counted plug DP at most
$(wh+1)\,K(\min(w,h))\le(wh+1)\,K(10)$ when a side is at most $10$. Both are at most
$(K(10)+300\,000)(wh+w+h+1)$, the bound of Section~3.5. Since a side at most $10$ gives $wh\le 10M$
and $K(10)>600\,000$, the count is also at most $(10M+1)\,K(10)$. This last step is arithmetic on
the two Lean bounds and is not itself a Lean theorem.

\subsection{Construction}

% lean: ZZN/Constr/Round.lean:420 solveR
% lean: ZZN/Constr/Solve2.lean:214 solveRec2
% lean: ZZN/Constr/RoundCount.lean:135 phi
% lean: ZZN/Constr/RoundCount.lean:143 phi_le
% lean: ZZN/Constr/RoundCount.lean:350 phi_del
% lean: ZZN/Constr/RoundCount.lean:357 phi_delT
% lean: ZZN/Constr/RoundCount.lean:90 w_le_of_none
% lean: ZZN/Timed/RoundTime.lean:695 solveRT_time
BlockLayer and PairLayer each carry a recursion counter, started at $w+h$, and each recursive call
deletes at least two lines. The number of \emph{open gaps} of an instance is the sum over both axes
of $1$ if a top block exists, $1$ if a bottom block exists, and the number of ordered pairs $(a,b)$
of endpoint coordinates that admit a middle block. It is at most $2\cdot(1+1+16)=36$, and every
block deletion lowers it, so BlockLayer deletes at most $36$ blocks. Without a block on an axis of
length $n\ge 13$, the least endpoint coordinate is at most $3$, the greatest at least $n-4$, and
neighbouring endpoint lines at most $10$ apart, so $n\le 3+3\cdot 10+4=37$. A round, one call of
BlockLayer, costs at most $80(M+1)^2+400$ steps outside its recursive call, and there are at most
$37$ rounds.

% lean: ZZN/Look/RoundTimeL.lean:48 leafLT_time
% lean: ZZN/Look/LeafT.lean:476 lookJT_time
% lean: ZZN/Look/RoundTimeL.lean:57 solveZT_thinL_time
% lean: ZZN/Look/RoundTimeL.lean:75 solveZT_boxL_time
% lean: ZZN/Timed/Solve.lean:221 solveZT_time
% lean: ZZN/Constr/Solve2.lean:272 leaf16
% lean: ZZN/Look/RoundTimeL.lean:97 solveZ3LT_main
The final round ends in one of two cases, and the bound adds the costs of both.
\begin{itemize}
\item An instance with a side at most $10$ from the start has area at most $10M$. Its final
  instance costs $(10M+1)\,K(10)$ for the plug DP and $3(10M+1)^2+10$ for checking the answer.
\item Otherwise both sides are at most $37$ after the blocks, and PairLayer runs. It makes at most
  $w+h\le 74$ recursive calls at $700\,(37+1)^2$ steps each, then the final solver, then the check
  of its answer on an area of at most $16\cdot 37$, at $3(16\cdot 37+1)^2+10$ steps. The counted
  final solver takes at most $(wh+1)\,K(\min(w,h))+28\,000\,000$ steps, of which the table lookup
  takes at most $28\,000\,000$. The bound covers every final instance with sides at most $37$ and
  a side at most $16$, so the final-solver term is $(10\cdot 37+1)\,K(10)+28\,000\,000$.
\end{itemize}
With the table lookup this gives the bound
\[
37\,(80(M+1)^2+400) + (10M+1)\,K(10) + 3(10M+1)^2 + 10 + B_L,
\]
\[
B_L = 74\cdot 700\cdot 38^2 + (10\cdot 37+1)\,K(10) + 28\,000\,000 + 3\,(16\cdot 37+1)^2 + 10 .
\]
% lean: ZZN/Timed/RoundTimeR.lean:44 solveZT_boxR_time
% lean: ZZN/Timed/RoundTimeR.lean:61 solveZ3RT_main
With the plug DP as final solver on every final instance, a side at most $16$ and the other at most
$37$, the final-solver term is $(16\cdot 37+1)\,K(16)$ and $B_L$ becomes
\[
B_R = 74\cdot 700\cdot 38^2 + (16\cdot 37+1)\,K(16) + 3\,(16\cdot 37+1)^2 + 10 .
\]
% values: paper/tools/kwr.py
With the values of $K(10)$ and $K(16)$ from \lean{paper/tools/kwr.py}, $B_L=350\,693\,907\,985$ and
$B_R=1\,514\,396\,612\,559\,921$. Expanded in $M$, the table-lookup bound has the coefficient
$3\,260$ of $M^2$, the coefficient $10\,K(10)+5\,980=9\,449\,872\,660$ of $M$, and the constant
term $351\,638\,912\,426$.

\section{Lean name map}\label{sec:names}

This section supports Section~4 by mapping the terms and statements of the paper to the Lean
declarations, with paths relative to \lean{proof/}. A statement is listed by what it says.

{\small
\begin{longtable}{@{}>{\raggedright\arraybackslash}p{0.32\textwidth}>{\raggedright\arraybackslash}p{0.34\textwidth}>{\raggedright\arraybackslash}p{0.28\textwidth}@{}}
\toprule
Paper & Lean & File \\
\midrule
\endfirsthead
\toprule
Paper & Lean & File \\
\midrule
\endhead
\multicolumn{3}{@{}l}{\emph{Definitions (Section 1.1)}} \\
adjacent, path & \lean{Adjacent}, \lean{IsPath} & \lf{ZZN/Statement.lean} \\
instance, well-formed & \lean{Inst}, \lean{Inst.WellFormed} & \lf{ZZN/Statement.lean} \\
solution, solvable & \lean{IsSolution}, \lean{Solvable} & \lf{ZZN/Statement.lean} \\
sign $\sigma$ & \lean{sgn} & \lf{ZZN/Statement.lean} \\
delete two lines, transpose & \lean{Inst.deleteAt}, \lean{Inst.transpose} & \lf{ZZN/Statement.lean} \\
insert two lines & \lean{Inst.insertAt} & \lf{ZZN/BandExtension.lean} \\
free line & \lean{Inst.FreeLine} & \lf{ZZN/Reductions.lean} \\
wide & \lean{InDom} & \lf{ZZN/Skeleton.lean} \\
IPS-admissible & \lean{HasHamPath} & \lf{GridHam/Basic.lean} \\
IPS characterization & \lean{ips\_characterization}, \lean{IsForbidden} & \lf{GridHam/Main.lean}, \lf{GridHam/Forbidden.lean} \\
\multicolumn{3}{@{}l}{\emph{Thin rectangles (Section 3.1)}} \\
plug DP & \lean{thinBF}, \lean{stepCell}, \lean{stageF} & \lf{ZZN/Statement.lean} \\
the plug DP accepts exactly the solvable well-formed instances & \lean{thinBF\_iff} & \lf{ZZN/Thin/Fast.lean} \\
plug DP with paths, which returns only solutions, and one for every well-formed solvable instance & \lean{thinSolveW2}, \lean{thinSolveW2\_sound}, \lean{thinSolveW2\_complete} & \lf{ZZN/Constr/ThinWit2.lean} \\
state bound $N(b)$ & \lean{Nw}, \lean{wcnt}, \lean{card\_fold} & \lf{ZZN/Timed/ThinRopeCost.lean}, \lf{ThinWord.lean}, \lf{ThinWordCount.lean} \\
\multicolumn{3}{@{}l}{\emph{Catalogue and necessity (Section 3.2)}} \\
fires, passes & \lean{fires3}, \lean{Passes} & \lf{ZZN/Statement.lean} \\
corner frames, frame entries & \lean{Frame}, \lean{frames}, \lean{frameFires} & \lf{ZZN/Statement.lean} \\
clockwise order & \lean{perimIndex} & \lf{ZZN/Statement.lean} \\
P, T1, T2 & \lean{parityOk}, \lean{t1Fires}, \lean{t2Fires} & \lf{ZZN/Statement.lean} \\
L1, L6 & \lean{l1Fires}, \lean{l6Fires} & \lf{ZZN/Statement.lean} \\
L2, L3, L4, L5 & \lean{testL2}, \lean{testL3}, \lean{testL4}, \lean{testL5} & \lf{ZZN/Statement.lean} \\
E1 & \lean{edgeClosure} & \lf{ZZN/Statement.lean} \\
B & \lean{boundaryOnly}, \lean{BOUNDARY3\_EVEN}, \lean{BOUNDARY3\_ODD} & \lf{ZZN/Statement.lean} \\
C, D & \lean{corner4}, \lean{CORNER4}, \lean{CORNER4\_ODD} & \lf{ZZN/Statement.lean} \\
R & \lean{effAlt3} & \lf{ZZN/Statement.lean} \\
R1, R2, R3 & the three cases of \lean{rewritesAt} & \lf{ZZN/Statement.lean} \\
settled cells, effective endpoints & \lean{RWState}, \lean{rwStep}, \lean{rewriteAll} & \lf{ZZN/Statement.lean} \\
bypass walk & \lean{walk3}, \lean{staircase} & \lf{ZZN/Statement.lean} \\
every solvable wide instance passes the catalogue & \lean{theoremB\_K} & \lf{WinKU/Facts.lean} \\
P does not fire on a solvable instance & \lean{parity\_of\_solvable} & \lf{ZZN/Parity.lean} \\
separation fact & \lean{parP\_chain} & \lf{ZZN/Planar.lean} \\
two disjoint chains in the grid do not have alternating ends on the border, T1 does not fire on a solvable well-formed instance with sides at least 2 & \lean{t1\_core}, \lean{t1\_sound} & \lf{ZZN/SoundT1.lean} \\
T2 does not fire on a solvable well-formed instance & \lean{t2\_sound} & \lf{ZZN/SoundT2.lean} \\
Lean statements behind the lemma on the local entries, for a solvable well-formed instance of any size: L1 and L6 do not fire (L6 for sides at least 2), and L2 to L5 and E1 do not fire in the identity frame & \lean{l1\_sound}, \lean{l6\_sound}, \lean{testL2\_sound} to \lean{testL5\_sound}, \lean{edgeClosure\_sound} & \lf{ZZN/SoundLocal.lean} \\
no frame entry fires in any of the 8 frames of a solvable wide instance, given the certificates & \lean{frames\_sound} & \lf{ZZN/TheoremB.lean} \\
forced parts of a path at its endpoints, the run between the effective endpoints & \lean{path\_trim}, \lean{endpoint\_info} & \lf{ZZN/SoundR.lean} \\
R does not fire on a solvable wide instance & \lean{r\_sound} & \lf{ZZN/SoundRMain.lean} \\
settled-cell shapes & \lean{remShape\_rewriteAll} & \lf{ZZN/PassTR.lean} \\
window certificate, completion test & \lean{certify}, \lean{outsideOK} & \lf{ZZN/WinCheck.lean} \\
an instance whose window configuration has a certificate has no solution & \lean{window\_sound} & \lf{ZZN/WinSound.lean} \\
list of window certificates & \lean{CertFacts}, \lean{certFactsK} & \lf{ZZN/WinBCD.lean}, \lf{WinKU/Facts.lean} \\
certificate transfer (Section~\ref{sec:transfer}) & \lean{bSound\_of}, \lean{cSound\_of} & \lf{ZZN/WinBCD.lean}, \lf{ZZN/C8R.lean} \\
a match of a C configuration marked $*$ makes R fire & \lean{c8\_firesF}, \lean{c8\_baseF}, \lean{c8\_sound} & \lf{ZZN/C8R.lean} \\
a match of a B triple marked $\dagger$ makes R fire & \lean{b5\_all}, \lean{b5\_stepR}, \lean{b5\_base} & \lf{ZZN/B5Fire.lean}, \lf{ZZN/B5R.lean}, \lf{ZZN/B5Base.lean} \\
\multicolumn{3}{@{}l}{\emph{Deleting and inserting free lines (Section 3.3)}} \\
edge strip, wide-gap compression & the first three cases of \lean{RedX} & \lf{ZZN/Skeleton.lean} \\
distant-endpoint compression & \lean{RedB}, \lean{CompCore}, and for $8\le d$, $d+10\le w$ also the fourth case of \lean{RedX} & \lf{ZZN/PassB.lean}, \lf{ZZN/PassX.lean}, \lf{ZZN/Skeleton.lean} \\
a pair deletion keeps an instance wide & \lean{inDom\_deleteAt}, \lean{inDom\_deleteB} & \lf{ZZN/Skeleton.lean}, \lf{ZZN/PassB.lean} \\
a pair deletion keeps a passing wide instance passing & \lean{passX\_all}, \lean{passB} & \lf{ZZN/PassB.lean} \\
transposition keeps a wide instance passing (Section~\ref{sec:rtrans}) & \lean{passT}, \lean{effAlt3\_transpose} & \lf{ZZN/PassTR.lean} \\
band extension & \lean{band\_extension}, \lean{bandExtE}, \lean{seed\_or\_flip} & \lf{ZZN/BandExtension.lean}, \lf{ZZN/Statement.lean} \\
a solution of the smaller instance of a pair deletion gives one of the larger & \lean{lift\_redX}, \lean{lift\_redB} & \lf{ZZN/Skeleton.lean}, \lf{ZZN/PassB.lean} \\
\multicolumn{3}{@{}l}{\emph{Sufficiency (Section 3.4)}} \\
a wide instance is solvable if and only if it passes the catalogue & \lean{passes\_iff\_solvable\_K} & \lf{FinalN.lean} \\
induction on the area & \lean{theoremA\_B}, \lean{theoremA\_boxes} & \lf{ZZN/PassB.lean}, \lf{ZZN/BoxK/Box.lean} \\
an axis of length at least 23 admits a pair deletion & \lean{redX\_of\_wide} & \lf{ZZN/Skeleton.lean} \\
an axis of length 17 to 22 admits a pair deletion & \lean{none\_all} & \lf{ZZN/BoxK/None.lean} \\
cut tree, checks & \lean{CTree}, \lean{check}, \lean{check\_sound} & \lf{ZZN/Cut/Tree.lean} \\
template, checks on an instance & \lean{TPos}, \lean{tcheck}, \lean{solvable\_of\_tcheck} & \lf{ZZN/Cut/Template.lean} \\
symmetries & \lean{Sym}, \lean{allSym} & \lf{ZZN/Moves.lean}, \lf{ZZN/Orbit.lean} \\
row tuple, column tuple & \lean{rowsOf}, \lean{colsOf} & \lf{ZZN/BoxK/Box.lean} \\
order of the tuples, canonical image & \lean{keyR}, \lean{key2}, \lean{Gbox} & \lf{ZZN/BoxK/Slice.lean}, \lf{ZZN/BoxK/Box.lean} \\
grid certificate, row class & \lean{BoxCert}, \lean{rowsOK}, \lean{SliceGood}, \lean{sliceCOK} & \lf{ZZN/BoxK/} \\
the 21 grid certificates check & \lean{boxes\_allK} & \lf{FinalN.lean} \\
a passing wide instance with $w\le h\le 22$ and no pair deletion is solvable (Lean also assumes that no same-colour split applies) & \lean{finite\_le} & \lf{ZZN/BoxK/Box.lean} \\
symmetries keep passing & \lean{passes\_lab}, \lean{passes\_transpose\_iff}, \lean{passes\_refl\_of\_image} & \lf{ZZN/Orbit.lean} \\
some reflection passes & \lean{passImgB} & \lf{ZZN/Finite.lean} \\
\multicolumn{3}{@{}l}{\emph{Algorithm (Section 3.5)}} \\
Decide, which returns true exactly on the solvable well-formed instances & \lean{solvableB}, \lean{solvableB\_iff} & \lf{ZZN/Statement.lean}, \lf{ZZN/Thin/Decide.lean} \\
SelectBlock & \lean{blkX}, \lean{leftK}, \lean{rightK}, \lean{innerK}, \lean{innerPair} & \lf{ZZN/Constr/Round.lean} \\
EdgeSplice, snake & \lean{edgeLift}, \lean{edgeLiftL}, \lean{colSnake} & \lf{ZZN/Constr/EdgeBlock.lean} \\
MiddleSplice & \lean{innerLift}, \lean{wideExt} & \lf{ZZN/Constr/Round.lean}, \lf{ZZN/Constr/WideBand.lean} \\
deleting a block is a sequence of pair deletions, and its lift maps solutions to solutions & \lean{blkX\_ok}, \lean{innerLift\_sol}, \lean{edgeLift\_sol}, \lean{edgeLiftL\_sol} & \lf{ZZN/Constr/Round.lean}, \lf{EdgeBlock.lean} \\
Construct, BlockLayer & \lean{solveZ3}, \lean{solveR} & \lf{ZZN/Constr/Round.lean} \\
PairLayer & \lean{solveZ2}, \lean{solveRec2} & \lf{ZZN/Constr/Solve2.lean} \\
final solver, table variant & \lean{leafL} & \lf{ZZN/Look/Complete.lean} \\
TableLookup, table & \lean{lookJ}, \lean{lookLe}, \lean{tabAll} & \lf{ZZN/Look/Leaf.lean}, \lf{ZZN/Look/Table.lean} \\
IPS recursion & \lean{ipsP}, \lean{ipsP\_valid} & \lf{ZZN/Look/IPS.lean} \\
Construct returns only solutions, and one for every well-formed solvable instance, table variant & \lean{solveZ3L\_main}, \lean{solveZ3\_completeL} & \lf{ZZN/Look/Main.lean}, \lf{ZZN/Look/Complete.lean} \\
the same, plug-DP variant & \lean{solveZ3\_sound}, \lean{solveZ3\_complete}, \lean{thinOK16} & \lf{ZZN/Constr/Round.lean}, \lf{ZZN/Timed/Construct.lean} \\
an axis of length at least 11 without a pair deletion has length at most 16 & \lean{leaf16} & \lf{ZZN/Constr/Solve2.lean} \\
step model & \lean{TimeM}, \lean{lmapT}, \lean{ldedupT} & \lf{ZZN/TimedStatement.lean} \\
counted Decide & \lean{solvableBRT}, \lean{solvableBRT\_iff} & \lf{ZZN/Timed/DecideR.lean} \\
counted Construct & \lean{solveZ3T}, \lean{solveZ3T\_ret} & \lf{ZZN/Timed/RoundT.lean} \\
$K_s(b)$, $K(b)$ & \lean{KsR}, \lean{KwR} & \lf{ZZN/Timed/ThinRopeCost.lean} \\
open gaps & \lean{phi}, \lean{phi\_le}, \lean{phi\_del}, \lean{w\_le\_of\_none} & \lf{ZZN/Constr/RoundCount.lean} \\
step bounds of the counted Decide, of the counted Construct with the table variant and with the plug-DP variant & \lean{solvableBRT\_main}, \lean{solveZ3LT\_main}, \lean{solveZ3RT\_main} & \lf{ZZN/Timed/DecideR.lean}, \lf{ZZN/Look/RoundTimeL.lean}, \lf{ZZN/Timed/RoundTimeR.lean} \\
\multicolumn{3}{@{}l}{\emph{Statements of the challenge files (Section 4)}} \\
counted Decide, counted catalogue on wide instances & \lean{solvableBRT\_main}, \lean{fires3T\_iff\_solvable} & \lf{ChallengeDecide.lean}, proved in \lf{ZZN/Timed/DecideR.lean}, \lf{ZZN/Timed/Main.lean} \\
counted Construct, plug-DP variant & \lean{solveZ3RT\_main} & \lf{ChallengeConstruct.lean}, proved in \lf{ZZN/Timed/RoundTimeR.lean} \\
\bottomrule
\end{longtable}
}

\bibliographystyle{plain}
\bibliography{references}

\end{document}
